% !TEX root = thesis.tex
% =========================================================================
% Chapter 3 — Related Work
%
% Simple-English pass over the prior draft. Same sections, same claims,
% same \cite keys, same structure, same table. Wording is simplified and
% sentences restructured; no information added, removed, or altered.
%
% Two grammatically broken sentences in the graph-IE section (PICK soft
% weights; SPADE serialiser-free) were made readable as part of the
% rewrite — meaning unchanged.
%
% LEFT AS-IS by request (option 1, not the full consistency pass):
%   - \textbf{} bold, the \S\ref{} cross-ref style, and the
%     ch:related-work label (note: other chapters use the chap: prefix).
%
% \ref FIXES (the old §?? boxes):
%   - sec:bgd-problem  (zero-shot, §2.4)  -> sec:bgd-zeroshot
%   - sec:bgd-problem  (LLM reasons, §2.2.1) -> removed; stated inline
%   - sec:bgd-evaluation-tuplef1 (§2.5)   -> chap:evaluation
%
% Cited keys — ALL already in references.bib; do NOT add anything:
%   long2021scene, nagayi2025foodocr, arief2026halalbench,
%   yazdani2025glinerbiomed, devlin2019bert, zaratiana2023gliner,
%   yu2020pick, hwang2021spade
% =========================================================================

\chapter{Related Work}\label{ch:related-work}

Chapter~\ref{chap:background} explained the ideas the pipeline relies on. This chapter sets those ideas against the published systems that the pipeline borrows from and moves away from. The comparison is then used to state the exact gap that the thesis fills. The review follows the same split the pipeline uses. \S\ref{sec:rw-ocr} looks at optical character recognition on food and nutrition labels, and at the evidence that recognition on its own is not enough for this kind of image. \S\ref{sec:rw-ie-paradigms} looks at the two supervision-free routes the classifier stage could have taken --- end-to-end extraction with large language and vision-language models, and zero-shot token classification --- and explains why the thesis limits a vision-language model to one sub-task instead of the whole problem. \S\ref{sec:rw-graph-ie} reviews graph-based information extraction, which is the closest neighbour of the association stage, and sets out exactly what the thesis takes from PICK and SPADE and what it does differently. \S\ref{sec:rw-zeroshot} places the work inside the zero-shot information-extraction literature. \S\ref{sec:rw-gap} draws the three positions together into the research gap and the choices that follow from it.

One thread runs through the whole comparison. Every system reviewed here needs task-specific training, gives up interpretability, or does both. This thesis aims at the one corner of the design space that is left: a modular, fully interpretable, zero-shot pipeline that rebuilds four-field nutrition tuples from any multilingual label. The sections below show that this corner is really empty.

% =========================================================================
\section{OCR on Food and Nutrition Labels}\label{sec:rw-ocr}
% =========================================================================

The detection-and-recognition cascade of \S\ref{sec:bgd-ocr} is well developed on the documents and scene text that fill most of the OCR literature~\cite{long2021scene}. Supplement and nutrition labels are a different and harder case, and the recent benchmarks aimed at them say this plainly. Nagayi et~al.~\cite{nagayi2025foodocr} test four open engines --- Tesseract, EasyOCR, PaddleOCR, and TrOCR --- on real food-packaging images, including nutrition-facts panels. They report two problems. Curved surfaces lower detection recall, and multilingual labels confuse language-specific recognition models. These are the same two failure modes that matter most for the dataset in this thesis. HalalBench~\cite{arief2026halalbench} is the first open multilingual benchmark for food-packaging OCR, and it reaches the same view at a larger scale. Across a fourteen-language corpus, even the strongest engines reach only a low F1, and all of them fail completely on scripts such as Japanese. A clustering-based post-processing step then recovers a large part of that loss.

Two things follow for this work. First, the evidence that OCR alone cannot produce clean structured nutrition data is exactly what motivates the semantic-classification and graph-association stages that form the contribution. The perceptual layer sets the ceiling (\S\ref{sec:bgd-ocr}), but that ceiling is too low to be useful on its own. Second, engine choice clearly matters, yet it is not the research question. The pipeline therefore takes the PP-OCRv5 engine off the shelf (\S\ref{sec:bgd-ocr}) and treats it as one component that can be swapped, rather than as a contribution. The benchmark literature both allows and supports this choice.

% =========================================================================
\section{Information Extraction Without Task Supervision}\label{sec:rw-ie-paradigms}
% =========================================================================
The classifier stage can run without any training on the target task. Two approaches make this possible, and the choice between them affects the rest of the design. The first treats the task as a single step: the label image is passed to a large model, and the model returns the tuples directly. The second keeps the stage as token classification, where each token is given a role by matching it to a short description of the target types rather than to labelled examples. The next two subsections take each one in turn, and show why the thesis follows the second for the main pipeline while limiting the first to a single sub-task.
\subsection{End-to-End Extraction with Language and Vision-Language Models}\label{sec:rw-llm}

The most direct alternative to the whole pipeline is to drop its stages and hand the label image to a large vision-language model, with an instruction to return tuples. This baseline is appealing. It needs no annotation, and in principle it copes with any layout. Even so, the thesis does not use it as the \emph{primary} method, for three reasons. Each one goes against a goal the thesis sets for itself. The first is \textbf{cost}. Calling a large model on every label at inference is expensive, and it scales poorly in a market where tens of thousands of new layouts appear each year. The second is the \textbf{absence of interpretability and transparency}. An end-to-end model shows no intermediate representation that can be inspected, so a wrong tuple cannot be traced to a stage, a token, and a decision --- which is exactly the kind of locatable failure the modular design is meant to give. The third is the \textbf{absence of an audit trail}. There is no per-image record that ties each output back to the evidence that produced it. This trade-off between cost and capability is real, not just a matter of taste, and the open-model literature shows it: GLiNER-BioMed~\cite{yazdani2025glinerbiomed} is presented as a practical alternative to large language models, which it describes as too costly and too large for resource-limited use. The thesis does not ban vision-language models. Instead, it \emph{confines} one to the association sub-task. The model is given an injected prompt, built from the outputs of the earlier stages, and is used only to bind tokens that have already been classified; the actual setup is given in the methodology chapter. Most of the pipeline therefore stays cheap and open to inspection, and the model is handed only the single relational step where its flexibility is worth the cost.

\subsection{Token Classification and Semantic Role Labelling}\label{sec:rw-classification}

In the terms of the literature, the classifier stage of \S\ref{sec:bgd-ner} is a token-classification or semantic-role-labelling problem: give each token a role in the target schema. Two learned approaches set the scene for the thesis's three classifier families. Pre-trained transformer encoders such as BERT~\cite{devlin2019bert} build a contextual vector for each token --- a vector that captures meaning in context, not in isolation. This helps rare or unclear tokens, because they can be classified from their surroundings rather than from their surface form alone. GLiNER~\cite{zaratiana2023gliner} builds on this idea. It treats named-entity recognition as matching candidate spans against natural-language type strings given at inference, so a new schema is reached by changing the prompt instead of retraining. Its domain-adapted version, GLiNER-BioMed~\cite{yazdani2025glinerbiomed}, takes this to specialised vocabularies. It reports state-of-the-art zero-shot results while staying far cheaper than a large language model, and it is the learned-NER model the thesis evaluates. These two sit next to the thesis's other two families (\S\ref{sec:bgd-ner}). One is a rule-based classifier, whose prior is stated openly. The other is an embedding-prototype classifier, whose prior is the distributional geometry of the BGE-M3 encoder (\S\ref{sec:bgd-embeddings}). All three aim at the same semantic-role-labelling goal. The embedding route, in particular, reaches it through cosine similarity to category prototypes, rather than through a trained span tagger. One cost is specific to a graph pipeline, and it shaped the design, so it is worth stating here. Learned NER works on serialised text, so its predicted spans have to be re-aligned to the OCR tokens and their bounding boxes before the graph stage can use them. This alignment is not trivial, and it brings its own failure modes when serialisation reorders or merges tokens. The rule-based and embedding routes never leave the token-and-geometry representation, so they do not pay this cost.

% =========================================================================
\section{Graph-Based Information Extraction}\label{sec:rw-graph-ie}
% =========================================================================

Association is about relations, not classification, and the natural structure for relational reasoning over a document is the graph from \S\ref{sec:bgd-graphs}. Two systems sit closest to the thesis's association stage. The contribution is clearest when it is stated as what it keeps from each and what it replaces. PICK~\cite{yu2020pick} treats key-information extraction as graph learning together with graph convolution. It combines text and visual features and learns a soft adjacency matrix end-to-end over the text segments of the document. The thesis takes PICK's central idea: adjacency should be continuous, not binary, so relations are held as soft weights rather than fixed yes-or-no edges (\S\ref{sec:bgd-graphs}). It departs from PICK on one point. It \emph{computes} the edge weights with closed-form functions of geometry, instead of learning them. This is what keeps the zero-shot commitment while still giving the soft-adjacency behaviour. SPADE~\cite{hwang2021spade} restates extraction as spatial-dependency parsing. It builds a directed relation graph over two-dimensional tokens; it needs no serialiser and runs end-to-end, going from OCR tokens and their coordinates to a dependency graph. SPADE also adds the edge-resolution step the thesis calls Tail-Collision Avoidance. The idea behind it is that a node in a semi-structured document almost always has a single incoming edge per relation. The thesis adopts this principle, but uses it as a light single-pass guard rather than the original iterative resolution; this is a deliberate simplification, and its limits are set out in the methodology chapter.

The clearest difference from both systems is supervision. PICK and SPADE are trained end-to-end on annotated corpora. Neither runs zero-shot, and there is no annotated corpus of supplement labels with tuple-level ground truth to train them on. The thesis therefore keeps the relational, soft-adjacency, TCA-guarded \emph{view} that these systems set out, while replacing their learned parts with closed-form, interpretable, training-free geometry. The exact edge-construction modes that make up the contribution --- proximity, structural rank-under-context, and their merged combination --- belong to the methodology chapter and are not shown here. This section claims only that the thesis inherits the graph formulation of PICK and SPADE, and breaks from them at the precise point where they need supervision.

% =========================================================================
\section{Zero-Shot Information Extraction}\label{sec:rw-zeroshot}
% =========================================================================

The working meaning of zero-shot used throughout this thesis was fixed in \S\ref{sec:bgd-zeroshot}: no part of the system is fitted on labelled examples of the target task, and the pipeline must still work on unseen labels, layouts, and languages. The literature above shows why this is the right framing for the domain, rather than a limit the thesis imposes on itself. There is no annotated corpus of supplement labels to train on. The set of products is open and keeps growing, so any model tied to a fixed set of templates would already be out of date when it shipped. The range of real labels is also very wide. HalalBench~\cite{arief2026halalbench} makes this concrete, with its fourteen languages and its finding that strong engines fail on unseen scripts --- direct evidence that a method which depends on per-layout or per-language training cannot scale. On the other side, the zero-shot named-entity-recognition line of GLiNER~\cite{zaratiana2023gliner} and GLiNER-BioMed~\cite{yazdani2025glinerbiomed} shows that high-quality extraction without fine-tuning is possible, and much cheaper than with a large model. This is why a learned open-vocabulary NER baseline is included in the classifier stage. What the literature does not settle is which kind of zero-shot prior suits this domain best. The three families of \S\ref{sec:bgd-ner} place their generalisation on an explicit prior, a transferred prior, and a distributional prior, in turn. Which one wins is an empirical question, answered in the experimental chapter rather than decided in advance by the literature.

% =========================================================================
\section{Research Gap}\label{sec:rw-gap}
% =========================================================================

The reviewed work marks out three positions, and the gap is best seen as the corner that none of them reaches. \textbf{OCR-only and regular-expression pipelines} break on the structural problems of supplement labels --- dense rows of numbers, multi-context columns, multilingual text, and broken layouts --- as the benchmark evidence of \S\ref{sec:rw-ocr} shows. \textbf{Graph-based extractors} such as PICK and SPADE rightly treat extraction as relational and use soft, geometry-aware adjacency, but they learn their structure end-to-end from annotated corpora and are not built for zero-shot use across unseen multilingual layouts. \textbf{End-to-end language and vision-language models} remove the need for training data, but they bring back the very costs the thesis avoids: expense at scale, an opaque mapping with no intermediate to inspect, and no audit trail.

The gap the thesis fills sits between these positions. It joins their strengths and avoids their costs. The first part is an explicit graph over OCR tokens that have been semantically classified in a zero-shot way. The novelty is in \emph{how} the edges and relations are built: by closed-form geometric and structural rules, not by learned weights. In this way the relational power of PICK and SPADE is kept, but without their supervision. The second part is an injected prompt that limits a vision-language model to the association sub-task alone, instead of using it as an end-to-end black box. The model's flexibility is then used only where it is worth the cost, and the rest of the pipeline stays interpretable. The output of this combination is the four-field tuple $\langle\text{nutrient},\,\text{quantity},\,\text{unit},\,\text{context}\rangle$, whose evaluation is defined in Chapter~\ref{chap:evaluation}; the reference base on which a value is stated is carried by the \emph{context} field.

In one sentence: \emph{no existing approach rebuilds four-field nutrition tuples from arbitrary, multilingual, previously unseen supplement labels without task-specific training while keeping every stage interpretable} --- and this is exactly what the thesis makes possible. It removes three concrete limits: the fragility of order- and regex-based pipelines on broken multilingual layouts, the reliance of graph-based extractors on annotated training corpora, and the opacity and cost of end-to-end model extraction. Table~\ref{tab:rw-positioning} sums up the positioning, and these three points lead straight into the methodology chapter.

\begin{table}[t]
  \centering
  \small
  \begin{tabular}{lcccc}
    \toprule
    Approach & Zero-shot & Interpretable & Graph-based & Multilingual \\
    \midrule
    OCR + regex / rules               & Yes     & Yes     & No  & Partial \\
    PICK \cite{yu2020pick}            & No      & Partial      & Yes & Partial \\
    SPADE \cite{hwang2021spade}       & No      & Partial & Yes & Partial \\
    End-to-end LLM / VLM              & Yes     & No      & No  & Yes     \\
    \textbf{This thesis}              & Yes     & interpretable and auditable     & Yes & Yes     \\
    \bottomrule
  \end{tabular}
  \caption{Positioning of this thesis against the reviewed approaches along the
  four properties that define the research gap. Each prior approach forfeits at
  least one; the contribution is to satisfy all four simultaneously.}
  \label{tab:rw-positioning}
\end{table}